\section{Methods}

\subsection{Dataset}

We use two dermatoscopic image datasets. The HAM10000 dataset \cite{Tschandl2018} contains 10,015 dermatoscopic images from diverse populations, annotated with seven diagnostic categories. For binary melanoma classification, we retain only melanocytic lesions: 7,818 samples comprising melanoma (\emph{mel}) and melanocytic nevus (\emph{nv}) classes, yielding a malignant--benign ratio of approximately 14:86. No skin tone labels are available in HAM10000.

The Dermatology Data Initiative (DDI) dataset provides 656 dermatoscopic images with diverse skin-tone annotations \cite{DDI2021}. We interpret its numeric skin-tone codes using the Fitzpatrick scale \cite{Fitzpatrick1988}: 12 (Light, Fitzpatrick I--II), 34 (Medium, Fitzpatrick III--IV), or 56 (Dark, Fitzpatrick V--VI). All three subgroups are evaluated; the Medium subgroup is reported as a secondary analysis because its smaller melanoma count makes its estimates less stable than the Light and Dark subgroups.

DDI is split into train/validation/test partitions using a 60/20/20 ratio, stratified by the cross-product of label $\times$ skin tone to preserve class balance across subgroups. The split is saved for reproducibility (see Appendix). The resulting partition sizes are: Train\,$=$\,393, Validation\,$=$\,131, Test\,$=$\,132. The test set contains 42 Light samples (10 melanoma), 48 Medium samples (8 melanoma), and 42 Dark samples (10 melanoma). \emph{Power note:} The Dark test subgroup contains $n=42$ images with only $\sim 10$ melanoma cases, yielding large bootstrap uncertainty (bootstrap SD $\sim 0.15$ AUROC and 95\% CI half-width $\sim 0.20$; see Table~\ref{tab:subgroup_auroc}). Single-split fairness comparisons are fragile by construction; between-seed variability (SD $\approx 0.145$--$0.152$ for Dark AUROC) exceeds all between-stage differences (paired delta $\approx +0.028$). Accordingly, 5-seed aggregation is primary and seed-42 is treated as one draw from a high-variance distribution. All experiments use the DDI dataset only ($n=656$) for fine-tuning and evaluation. No external validation is performed on independent datasets (e.g., Fitzpatrick17k, HAM10000 held-out test sets, or clinical multi-centre cohorts); generalizability has not been demonstrated and requires future work. Split disjointness is enforced via assertion checks on filenames, and a byte-level SHA256 overlap check confirms that no generated image shares bytes with any held-out image. Dataset composition is summarized in Table~\ref{tab:dataset}.

\subsection{Model Architecture}

We use EfficientNet-B0 \cite{Tan2019EfficientNet} as the backbone architecture, with the final classifier layer replaced by a single-output linear layer for binary melanoma classification. The full model contains 4,008,829 parameters. Weights are initialized from ImageNet pretraining. Input images are resized to $224 \times 224$ and normalized with ImageNet channel statistics (mean $[0.485, 0.456, 0.406]$, std $[0.229, 0.224, 0.225]$).

Transfer learning follows a partial-freezing strategy. During the fine-tuning and generated-augmentation stages, early convolutional feature extractor layers are frozen, and only the final two feature blocks (\texttt{features.6}, \texttt{features.7}) and the classifier head are trained, leaving 2,744,861 trainable parameters. This reduces the number of trainable parameters relative to full fine-tuning while allowing the model to adapt high-level representations to the target domain. The generated-image ablation study uses this identical frozen-layer configuration, so its arms differ only in the number of generated images.

\subsection{Training Pipeline}

The pipeline consists of three sequential stages, each initializing from the preceding stage's checkpoint. Stage 1 uses its own hyperparameters as a pretrained initialization; Stages 2 and 3 share learning rate, weight decay, batch size, frozen-layer configuration, and epoch budget, so the only intended difference between them is the addition of generated images. We report the one remaining difference---the epoch budget of the ablation arms---explicitly rather than describing the comparison as fully controlled. Hyperparameters are summarized in Table~\ref{tab:training_config}.

\paragraph{Stage 1: Baseline.}
The model is trained on the full HAM10000 mel/nv subset (80/20 stratified split) for up to 10 epochs with early stopping (patience\,$=$\,3). The binary cross-entropy loss \texttt{BCEWithLogitsLoss} is used with no class weighting. The optimizer is AdamW with learning rate $3 \times 10^{-4}$ and weight decay $10^{-4}$. A \texttt{ReduceLROnPlateau} scheduler (factor\,$=$\,0.5, patience\,$=$\,2) monitors validation AUROC. This stage uses different hyperparameters from Stages 2--3, which is acceptable as it serves only as a pretrained initialization.

\paragraph{Stage 2: Fine-tuning.}
The baseline checkpoint is loaded, feature extractor layers are frozen as described above, and the model is fine-tuned on the DDI training partition for up to 10 epochs with early stopping (patience\,$=$\,3). The malignant--benign class imbalance in DDI is addressed by computing \texttt{pos\_weight} dynamically from the training split as $n_{\text{neg}} / n_{\text{pos}}$. The optimizer is AdamW with learning rate $5 \times 10^{-5}$ and weight decay $10^{-4}$, applied only to unfrozen parameters. Batch size is 32.

\paragraph{Stage 3: Generated Augmentation.}
The fine-tuned checkpoint is loaded with the same frozen-layer configuration. The DDI training set is augmented with generated dark-skin melanoma images (described in Section~\ref{sec:synthetic_aug}) and trained for up to 10 epochs (patience\,$=$\,3). Learning rate $5 \times 10^{-5}$, weight decay $10^{-4}$, batch size 32, and early-stopping patience are identical to Stage 2.

\textbf{Stage 3 confound note:} Generating only dark-skin melanomas (29$\to$290) while keeping \texttt{pos\_weight}$= n_{\text{neg}}/n_{\text{pos}}$ computed on the pre-augmentation training split ($291/102=2.8529$) changes the combined training class ratio from 102/291 (pos/neg) to 392/291; the balanced ratio for the combined set is $291/392=0.7423$. This over-weights malignancy in the augmented dark-skin portion by a factor of $\approx 3.8\times$ relative to the true combined distribution. Consequently, the headline Stage 2$\to$3 contrast does \emph{not} isolate the effect of augmentation alone; attribution is confounded by class composition and weighting. We quantify this via the arm study (Section~\ref{sec:pw\_arms}).

\paragraph{Class-imbalance weighting and the Stage~2/Stage~3 contrast.}
The \texttt{pos\_weight} applied in Stage 3 is the Stage-2 value $n_{\text{neg}} / n_{\text{pos}} = 291 / 102 = 2.8529$, computed on the pre-augmentation training split. Adding 290 generated images raises the effective positive count to 392 against 291 negatives, so the balanced ratio for the set actually presented to the loss is $291/392 = 0.7423$. The retained weighting therefore over-weights the minority class relative to the augmented set by a factor of $\approx 3.8$. This is not an incidental detail: it is a class-imbalance change that accompanies the augmentation, and it confounds attribution of any Stage~2 to Stage~3 difference to the generated images alone. We report it explicitly rather than describing the contrast as a clean single-variable comparison, and Section~\ref{sec:pw_arms} isolates its effect by re-running Stage 3 under three weighting regimes from the same Stage-2 checkpoint. The headline result in Section~\ref{sec:results} is the pre-existing configuration, and we note it is not the configuration that isolates the intended effect.

\paragraph{Hyperparameter standardization.}
Stages 2 and 3 share learning rate, weight decay, batch size, frozen layers, and epoch budget, and differ in the added generated images and in the implied class ratio discussed above. The ablation arms (Section~\ref{sec:ablation}) hold all Stage-3 hyperparameters fixed and vary only the number of generated images. Stage 1 differs in learning rate and batch size but serves only as a pretrained initialization.

\subsection{Generated-Image Augmentation}
\label{sec:synthetic_aug}

Our generated images are \emph{saved photometric and geometric transforms}, not samples from a generative model: we use the Albumentations library \cite{Buslaev2020} to apply fixed random transforms to existing training lesions, so no GAN or diffusion model is involved. From the dark-skin melanoma examples in the DDI \emph{training split} ($n\,{=}\,29$), each source image is transformed 10 times, yielding 290 generated images. We refer to these as \emph{generated} rather than \emph{synthetic} throughout, to avoid implying a generative model. Generated images derive exclusively from training-split data and never from validation or test images.

Because these are transformed copies of individual lesions rather than independent samples, we verified their provenance rather than assuming it. For each of the 290 generated images we computed a 32-dimensional descriptor and searched for its nearest neighbour among DDI images under all 8 dihedral transforms; the 32 darkest principal components were used to suppress colour and intensity, which dominate raw-pixel distances. All 290 images matched a \emph{training}-split dark melanoma as the nearest neighbour, none matched a held-out image, and all 29 source lesions were exercised. Median similarity to the matched training source was 0.9985 and the median margin over the best held-out decoy was $+0.2706$ (minimum margin $+0.0766$). Exact-duplicate exclusion is handled separately by a byte-level SHA256 check. The full per-image provenance report is provided in the Appendix.

The augmentation pipeline applies the following transforms in sequence:
\begin{enumerate}
    \item \textbf{RandomBrightnessContrast}: brightness and contrast variation ($\pm 0.2$ range, probability\,$=$\,0.8)
    \item \textbf{HueSaturationValue}: hue shift ($\pm 10$), saturation shift ($\pm 20$), value shift ($\pm 15$), probability\,$=$\,0.7
    \item \textbf{GaussNoise}: additive Gaussian noise, probability\,$=$\,0.2
    \item \textbf{ElasticTransform}: smooth geometric deformation ($\alpha\,{=}\,0.5$, $\sigma\,{=}\,5$), probability\,$=$\,0.1
    \item \textbf{RandomRotate90}: 90$^\circ$ rotation multiples, probability\,$=$\,0.5
    \item \textbf{HorizontalFlip}: mirror reflection, probability\,$=$\,0.5
\end{enumerate}

All outputs are resized to $224 \times 224$ and saved as JPEG files. Generated images are added to the DDI training set via \texttt{ConcatDataset}, increasing the effective training size from 393 to 683 samples. They are never added to validation or test partitions. The transform pipeline is seeded, so the generated set is reproducible for a given seed. Split disjointness and byte-level SHA256 overlap with held-out images are verified as described in the Appendix.

\textbf{Augmentation diversity (limitation):} The 290 generated images are photometric and geometric transforms of only 29 source lesions (10 variants each). No FID score is computed, no clinician review of lesion preservation was performed, and no comparison to generative models (e.g., diffusion/GAN) is made. Provenance checks show median descriptor similarity 0.9985 to the matched training source (confirming near-duplicates of the originals), with the minimum margin over the best held-out decoy $+0.0766$ and byte-level SHA256 excluding exact duplicates. Thus diversity is bounded by the source distribution.

\subsection{Class-Imbalance Arm Study}
\label{sec:pw_arms}

To separate the effect of augmentation from the accompanying change in class ratio, we re-run Stage 3 from the \emph{same} Stage-2 checkpoint under three arms that differ only in \texttt{pos\_weight} and generated-image composition. All arms share seed, split, frozen layers, learning rate, weight decay, batch size, and epoch budget.

\begin{itemize}
    \item \textbf{static}: \texttt{pos\_weight} $= 2.8529$ (the pre-augmentation ratio), 290 dark-melanoma generated images. This is the configuration underlying the headline result.
    \item \textbf{recomputed}: \texttt{pos\_weight} $= 0.7423$, the ratio of the actual combined set, with the same 290 generated images. This isolates the weighting change from any change in generated volume.
    \item \textbf{symmetric}: 101 dark-benign generated images added (spread over the 95 dark benign training lesions) so the combined set is exactly balanced at 392 positives and 392 negatives, with \texttt{pos\_weight} $= 1.0$. Here augmentation changes neither the class ratio nor the loss weighting.
\end{itemize}

This design isolates how much of the Stage-2 to Stage-3 change is attributable to the generated images rather than to the reweighting they induce.

\subsection{Evaluation}

The primary metric is the Area Under the Receiver Operating Characteristic curve (AUROC) \cite{Hanley1982}, which is threshold-independent and thus suitable for comparing models at different operating points. Secondary threshold-dependent metrics---sensitivity, specificity, positive predictive value (PPV), and F1 score---are also reported.

Classification thresholds for all three DDI test evaluations use Youden's $J$ statistic (maximizing $\text{TPR} - \text{FPR}$) computed on the corresponding validation predictions. AUROC is the primary evaluation metric as it is threshold-independent and appropriate for cross-stage comparison. Sensitivity, specificity, PPV, and F1 are threshold-dependent and computed using per-stage Youden's $J$ thresholds optimized on the corresponding validation set; as such, they are not directly comparable across stages. No matched-threshold analysis (fixed operating point across stages) was performed. For example, Light versus Dark sensitivity values reflect different operating points chosen per stage/validation split. The tables that report threshold-dependent metrics carry a note to that effect.

For uncertainty quantification, bootstrap 95\% confidence intervals are computed using 1{,}000 resampling iterations with the percentile method \cite{Efron1986}. Iterations that produce degenerate samples (single-class bootstrap draws) are discarded. Bootstrap analysis uses a fixed random seed (seed\,$=$\,0) for the random number generator. A permutation-based p-value is computed for the Dark AUROC improvement between Stages 2 and 3: the proportion of bootstrap AUROC differences $\leq 0$ across 1{,}000 paired resamples.

Subgroup analysis stratifies predictions by Fitzpatrick skin type code: Light (code 12, Fitzpatrick I--II), Medium (code 34, Fitzpatrick III--IV), and Dark (code 56, Fitzpatrick V--VI). All three are reported. Because the test set contains 8 melanomas in the Medium subgroup versus 10 in each of Light and Dark, Medium estimates carry wider uncertainty and are interpreted as secondary. Given the small subgroup sizes ($n\,{=}\,42$--$48$, 8--10 melanomas), all subgroup conclusions should be treated as preliminary and directional.

\subsection{Five-Seed Robustness}

Single-split results on 132 test images can move substantially under resampling, so we repeat the entire three-stage pipeline across 5 seeds (42--46). For each seed we regenerate the stratified DDI split, retrain the Stage-1 HAM10000 baseline, retrain Stages 2 and 3, regenerate the 290 generated images with that seed's transform seed, and re-derive all metrics on that seed's test set. We report the mean, sample standard deviation (ddof\,$=$\,1), and a seed-level percentile bootstrap 95\% confidence interval over the 5 seeds for each reported quantity. The Stage-1 baseline is retrained per seed rather than reused, so the reported variation reflects the full pipeline.
